\section{A polynomial sample lower bound at each fixed critical depth}
\label{sec:critical-polynomial-sample-lower}

A low-rate interpolating halfspace gives a sample lower bound for the
actual critical posterior. The bound is polynomial in the parity and
ambient dimensions, and the width limit is taken before those dimensions
grow.

\begin{lemma}[Amplifying a first-feature mean through its second moment]
\label{lem:critical-mixture-mean-seed}
Let $\mu_1$ be the uncontrolled first-feature law. Suppose a feature law
$\nu$ satisfies $D(\nu\Vert\mu_1)\le D$, where $D\ge1$, and its mean
$h$ satisfies $g(x)h(x)\ge a$ at every input, with $0<a\le1$ and
$g$ Boolean. Set $R=D/a^2$ and $q=R^{1/L}$. If $q\le D$, then
\[
 J_L(g)\le5000L R^{1/L}.
\]
\end{lemma}

\begin{proof}
Put $\theta=q/D\le1$. The first controlled law
$(1-\theta)\mu_1+\theta\nu$ has entropy at most $\theta D=q$ by
convexity, and its second moment dominates
$\theta hh^T$. Thus it has a tracked rank-one vector
$b_1=\sqrt\theta h$ with
\[
 g(x)b_1(x)\ge\sqrt\theta a=q^{-(L-1)/2}.
\]
This is a second-moment component; it need not equal the first
controlled law's mean.

For $L=2$, choose the readout $z=h/(2a)$. Its cost is at most
$1/(8\theta a^2)=q/8$, proving the assertion in this case.
Suppose $L\ge3$ and put $c=1/100$. For every $0\le v\le1$ and
$u\ge0$, the elementary estimate used in the constant-target
construction gives
\[
 \mathbb E\tanh(u+\sqrt v Z)\ge c\min\{u,1\}.
\]
At the second hidden layer shift the reference Gaussian by
$\sqrt q\,b_1$. Its mean-shift entropy is at most $q/2$, because
$Q_1\succeq b_1b_1^T$. After tanh, its actual mean $b_2$ satisfies
$g(x)b_2(x)\ge c q^{-(L-2)/2}$. Its second moment dominates
$b_2b_2^T$.

At each hidden layer $s=3,\ldots,L-1$, shift the reference Gaussian
by $(\sqrt q/c)b_{s-1}$. The entropy is at most $q/(2c^2)$, and
the new feature mean satisfies
\[
 g(x)b_s(x)\ge c q^{-(L-s)/2}.
\]
The induction uses only that every reference diagonal is at most one,
not that the diagonals or feature magnitudes are equal. If some shift
coordinate exceeds one in magnitude, the displayed minimum estimate
still gives the required lower bound. All shift vectors lie in the
corresponding covariance range; the formulas therefore also apply to
singular covariances. Tanh pushforward cannot increase the entropy.

Finally take $z=b_{L-1}/(2c q^{-1/2})$. It has the required global
margin and Gaussian readout cost at most $q/(8c^2)$. Summing all
costs gives the claimed bound.
\end{proof}

\begin{lemma}[A low-rate majority interpolant for parity samples]
\label{lem:critical-oriented-majority}
Let the labels be a parity on a nonempty coordinate set $S$, and draw
$n\ge1$ independent uniform cube inputs. Fix one $i_0\in S$ and an
odd number $m\le d-1$ of other coordinate indices $I$.
There is a sample-dependent Boolean majority function
\[
 g(x)=\operatorname{sign}\!\left(\sum_{j\in I}s_jx_j\right),
 \qquad s_j\in\{-1,1\},
\]
which fits every training label with probability at least
$1-n\exp[-m/(6n)]$.
For every sample, this function admits a first-feature law $\nu$ with
\[
 D(\nu\Vert\mu_1)\le dm,
 \qquad g(x)\mathbb E_\nu H(x)\ge1/100
       \quad\text{for every input }x.
\]
Consequently, whenever $dm\ge10000$,
\[
 J_L(g)\le5000L(10000dm)^{1/L}
       \qquad(L\ge2).
\]
\end{lemma}

\begin{proof}
Write $Y_i$ for the parity label of training row $i$, and choose
\[
 s_j=\operatorname{sign}\!\left(\sum_{i=1}^nY_iX_{ij}\right),
\]
with either fixed convention for a zero sum.
Conditional on the entire label vector $Y$, all coordinate columns
other than $i_0$ are independent uniform sign vectors: the missing
parity coordinate is uniquely determined by these columns and $Y$.
Thus $A_{ij}=Y_iX_{ij}$, for $j\in I$, are independent uniform sign
columns conditionally on $Y$.

If $T=\sum_{i=1}^n\varepsilon_i$ is a sum of independent uniform
signs, exchangeability gives
\[
 \mathbb E\left[\varepsilon_i\operatorname{sign}(T)\right]
       =\frac{\mathbb E|T|}{n}\ge\frac1{\sqrt{3n}}.
\]
Indeed $\mathbb ET^2=n$ and $\mathbb ET^4\le3n^2$, while Holder's
inequality gives
$\mathbb ET^2\le(\mathbb E|T|)^{2/3}(\mathbb ET^4)^{1/3}$.
For each training row, the variables $s_jA_{ij}$ are therefore
independent across $j$, bounded by one in absolute value, and have
mean at least $1/\sqrt{3n}$. Hoeffding's inequality bounds the
probability that their sum is nonpositive by $e^{-m/(6n)}$.
A union bound proves the training-fit probability. Dependence between
different training rows is irrelevant to that union bound.

Since $m$ is odd, the affine form defining $g$ is a nonzero integer
at every cube input, so its global signed margin is at least one.
Shift the first-row Gaussian weight mean to have coefficient $s_j$
on $I$ and zero elsewhere, including its bias. This costs exactly
$dm/2$ in relative entropy. At each input its residual Gaussian
variance is $(d+1)/d\le2$, and its feature mean is
\[
 \mathbb E\tanh\!\left(\sum_{j\in I}s_jx_j+
                       \sqrt{(d+1)/d}\,Z\right).
\]
The elementary mean-response bound from the constant-target
construction gives signed mean at least $1/100$ everywhere.
Data processing therefore proves the first-feature-law assertion.
Apply Lemma~\ref{lem:critical-mixture-mean-seed} with
$D=dm$ and $a=1/100$. Its hypothesis
$(10000D)^{1/L}\le D$ holds whenever $D\ge10000$ and $L\ge2$.
\end{proof}


\begin{theorem}[Critical parity sample lower bound]
\label{thm:critical-polynomial-sample-lower}
Fix $L\ge2$ and $0<\eta<1$. Let the target be parity on $k\le d$
specified input coordinates, and put
\[
 M=\min\{k,d^{1/(L-1)}\}.
\]
Suppose $M\to\infty$ and $1\le n\le M^{1-\eta}$. Draw $n$ independent
uniform cube inputs labeled by this parity. For all sufficiently large
$M$, with training-sample probability at least $1-1/(2M)$, the globally
margin-conditioned posterior at
$\gamma=1$ has Gibbs error tending to $1/2$ as $N\to\infty$ at
this fixed $d,L,n$. If $p_N(x)$ is its posterior probability of the
correct label, then on the same event
\[
 \liminf_{N\to\infty}\Pr_x\{p_N(x)<2/3\}\ge1/4.
\]
For full parity, this gives $n\le d^{(1-\eta)/(L-1)}$. If
$d\ge k^{L-1}$, it gives $n\le k^{1-\eta}$.
\end{theorem}

\begin{proof}
Let $m$ be the smallest odd integer at least
$12n\log(2M(n+1))$. Since $n\le M^{1-\eta}$, one has
$m=O(n\log(M+1))=o(M)\le o(d)$, so $m\le d-1$ for large $M$.
Lemma~\ref{lem:critical-oriented-majority} supplies, with probability at
least $1-1/(2M)$, a training-consistent halfspace and a first-feature
law of entropy at most $D=dm$ with mean margin $a=1/100$.

Apply Lemma~\ref{lem:critical-mixture-mean-seed}. Its hypothesis
$(D/a^2)^{1/L}\le D$ is equivalent to $D^{L-1}\ge10000$ and holds
for large $M$. The resulting training-consistent competitor has rate
at most
\[
 5000L(10000dm)^{1/L}.
\]

Any Boolean function depending on every one of the $k$ parity coordinates
has at least $k$ relevant augmented coordinates. By
Theorem~\ref{thm:critical-fixed-depth-ordering}, its critical rate is at least
\[
 \frac1{64}\min\{d^{1/(L-1)},(dk)^{1/L}\}
       =\frac1{64}(dM)^{1/L}.
\]
The ratio of the
preceding competitor upper bound to this lower bound is at most
\[
 C L\,M^{-\eta/L}[\log(M+1)]^{1/L},
\]
which tends to zero. Thus there is a positive rate gap between a
training-consistent competitor and every function retaining all
parity coordinates.

The exact finite-domain posterior-rate principle now implies that
posterior mass on the latter functions tends to zero in the width
limit. A function omitting any one parity coordinate has uniform
parity error exactly $1/2$, even if it depends on irrelevant coordinates.
This proves the Gibbs limit. Finally
$\mathbb E_xp_N(x)\to1/2$ and Markov's inequality give
$\Pr_x\{p_N(x)<2/3\}\ge1-(3/2)\mathbb E_xp_N(x)$.
\end{proof}

The theorem concerns the actual posterior, but takes width to infinity
at fixed dimension, depth, and sample size. It gives no finite-width
threshold uniform in growing dimension and no corresponding lower
bound for posterior-majority prediction at confidence threshold $1/2$.

\begin{corollary}[Slowly growing learner depth]
\label{cor:critical-growing-depth-sample-lower}
Fix $0<\eta<1$. For full parity, the same conclusion holds at sample sizes
$n\le d^{(1-\eta)/(L-1)}$ along integer depths $L=L(d)\ge2$ satisfying
\[
 L=o\!\left(\sqrt{\frac{\log d}{\log\log d}}\right).
\]
For each $d$, the width limit is taken with its corresponding $L(d)$
and sample size fixed.
\end{corollary}

\begin{proof}
Here $M=d^{1/(L-1)}\to\infty$. The halfspace construction has
$m\le Cn\log(M+1)=o(M)\le o(d)$ uniformly in $L\ge2$ at the
specified sample sizes. Its first-feature entropy bound and amplification
are also uniform. The rate ratio is therefore at most
\[
 C L\,d^{-\eta/[L(L-1)]}[\log(d+1)]^{1/L}
\]
with universal $C$. Its logarithm tends to minus infinity under the
stated depth assumption. The mixture hypothesis and the training
success probability also hold uniformly, so the same rate-gap argument
applies.
\end{proof}
